\documentclass[exercice]{thedocument}%
\startdatakeys
\source{}
\BO{2026}
\cycle{4}
\competencesDuSocle{CA,MO}
\objectifsApprentissage{C4D1aa,C4D1ab,C4D1ac}
\automatismes{C4D1ab,C4D3af}
\enddatakeys
\begin{document}%
Sarah se rend au marché une fois par semaine. Les quatre premières semaines,
elle achète successivement:
\begin{itemize}
    \item 500 g de fraises pour 4 euros
    \item 600 g de fraises pour 4,80 euros
    \item 450 g de fraises 3,60 euros
    \item 800 g de fraises pour 5,20 euros
\end{itemize}
\begin{enumerate}
    \item Présenter ces données dans un tableau.
    \item Ce tableau traduit-il une situation de proportionnalité~?
\end{enumerate}
\end{document}
